\chapter{Congruences}\label{ch:congruences}

\begin{center}\itshape
What are the last two digits of $3^{3^{3^3}}$, a number with more than $3\times10^{12}$ digits?\\
Arithmetic modulo $n$ answers in four lines.
\end{center}


On a twelve-hour clock, $14$ and $2$ name the same position. Their difference, $12$, is a whole number of turns. Congruences make that idea exact: $a\equiv b\pmod m$ means $m$ divides $a-b$. The sign $\equiv$ is not approximate equality. It records equality of residue classes.

Addition and multiplication preserve this relation. Division needs care. Modulo $6$ we have $3\cdot1\equiv3\cdot3$, but $1\not\equiv3$. The factor $3$ cannot be cancelled because it has no inverse modulo $6$. Modulo $7$ it can: $3\cdot5\equiv1$, so multiplication by $5$ undoes multiplication by $3$. A small example exposes the exact hypothesis that keeps algebra legal.

Now imagine two clocks ticking at once, one modulo $3$, one modulo $5$. A remainder on each clock determines one remainder modulo $15$. The Chinese remainder theorem is this reconstruction, with Bézout's identity supplying the formula. When the moduli share a factor, the remainders must agree where their information overlaps.

Congruences will also tell us about roots of equations modulo prime powers. A root modulo $p$ may continue to roots modulo $p^2,p^3,\ldots$, or it may not. Before trusting a large modular computation, we will learn which steps preserve solutions and what data are needed to lift one.


\section{The ring \texorpdfstring{$\Z/n\Z$}{Z/nZ}}

% t1-pass2 motive b31b4a00
A remainder is not a number with a permanent label. The integers $2$, $7$ and $-3$ represent the same class modulo $5$. Before we add or multiply such classes, we must prove that the answer does not depend on which representative we choose. That is what makes $\Z/n\Z$ a ring, rather than just a convenient collection of clock positions. Dividing is a separate question: an element can be cancelled exactly when it has a multiplicative inverse.

For $n\ge1$, $a\equiv b\pmod n$ means $n\mid a-b$. This is an equivalence relation whose classes are $\bar a=a+n\Z$, and there are exactly $n$ of them, $\bar0,\dots,\overline{n-1}$ (Euclidean division gives each $a$ a unique remainder).

\begin{figure}[ht]
\centering
\fig{ch03-clock}
\caption{The clock $\Z/12\Z$. Adding $5$ is a step of five hours. Because $\gcd(5,12)=1$, repeated addition visits every residue exactly once before returning. Dark grey rings mark the residues coprime to $12$; rust dots mark the others. Adding $4$, which shares a factor with $12$, stays inside $\{0,4,8\}$.}
\label{fig:clock}
\end{figure}

\subsection*{A class is a whole family, a representative is one member}
Modulo $5$, the integers $-3$, $2$, $7$ and $12$ belong to the same residue class: each differs from $2$ by a multiple of $5$. The class is the infinite set $2+5\Z$, not just whichever representative we choose to write. That is why $\bar2=\bar7$ is an equality of \emph{classes}. When we multiply classes, we have to prove that choosing $2$ or $7$ makes no difference. The next proposition checks the arithmetic behind that familiar habit.

A \emph{commutative ring with identity} is a set with addition and multiplication obeying the familiar associative, commutative and distributive laws, with identity elements satisfying $x+0=x$ and $x\cdot1=x$, and with an additive inverse $-x$ satisfying $x+(-x)=0$ for every $x$. We do \emph{not} demand multiplicative inverses. An element $x$ is a \emph{unit} if some $u$ satisfies $xu=1$. A \emph{field} is a commutative ring with $0\ne1$ in which every non-zero element is a unit. These definitions distinguish the two algebraic worlds we need, without assuming an earlier algebra course.

\begin{proposition}\label{prop:ringZn}
If $a\equiv a'$ and $b\equiv b'\pmod n$, then $a+b\equiv a'+b'$ and $ab\equiv a'b'\pmod n$. Hence $\bar a+\bar b=\overline{a+b}$ and $\bar a\bar b=\overline{ab}$ are well defined and make $\Z/n\Z$ a commutative ring with $n$ elements.
\end{proposition}
\begin{proof}
Suppose $a\equiv a'$ and $b\equiv b'\pmod n$.
Then $n$ divides both $a-a'$ and $b-b'$.
For addition,
\[
(a+b)-(a'+b')=(a-a')+(b-b'),
\]
a sum of multiples of $n$.
For multiplication,
\[
ab-a'b'=a(b-b')+b'(a-a'),
\]
also a multiple of $n$.
Thus changing representatives changes neither
the sum class nor the product class.

This proves that the operations on residue classes
are well defined. Their associative, commutative
and distributive laws follow by choosing integer
representatives and using the same laws in $\Z$.
The classes of $0$ and $1$ act as the additive
and multiplicative identities, and the class of
$-a$ is the additive inverse of the class of $a$.
Euclidean division gives each class exactly one
representative in $\{0,\ldots,n-1\}$, so there
are exactly $n$ classes.
\end{proof}

\begin{proposition}\label{prop:units}
$\bar a$ is invertible in $\Z/n\Z$ if and only if $\gcd(a,n)=1$. Hence $\Z/n\Z$ is a field if and only if $n$ is prime, and the group of units $(\Z/n\Z)^\times$ has $\varphi(n)$ elements.
\end{proposition}
\begin{proof}
$\bar a\bar u=\bar1$ means $au+nv=1$ for some $v$, which is B\'ezout's characterisation of $\gcd(a,n)=1$. The inverse is computed by the extended Euclidean algorithm. So the units are the classes of the integers in $\{1,\dots,n\}$ prime to $n$, which is the definition of $\varphi(n)$. Finally $\Z/n\Z$ is a field iff every $a\in\{1,\dots,n-1\}$ is prime to $n$ (and $n\ge2$), iff $n$ is prime.
\end{proof}

% t1-pass2 why
\paragraph{Why this works.}
An inverse modulo $n$ is not a number obtained by ordinary division. It is a class $\bar u$ with $\bar a\bar u=\bar1$, or equivalently an integer equation $au+nv=1$. B\'ezout's theorem says that equation has a solution exactly when $\gcd(a,n)=1$. Thus the nonzero classes all have inverses precisely for a prime modulus. This is why a field behaves differently from a general ring of remainders.

The units form a \emph{group under multiplication}: their product is a unit, multiplication is associative, $1$ is a unit, and every unit has a multiplicative inverse among the units. For example $(xy)^{-1}=y^{-1}x^{-1}$, since multiplying the two expressions gives $1$. This group is finite here, and its size is $\varphi(n)$. Euler's proof below uses only the fact that multiplication by a unit permutes this finite set.

\begin{theorem}[Euler, Fermat]\label{thm:eulerthm}
If $\gcd(a,n)=1$ then $a^{\varphi(n)}\equiv1\pmod n$. In particular, for $p$ prime and $p\nmid a$, $a^{p-1}\equiv1\pmod p$.
\end{theorem}
\begin{proof}
Multiplication by $\bar a$ is a bijection of the finite group $(\Z/n\Z)^\times$ (its inverse is multiplication by $\bar a^{-1}$). So the product of all units $x$ equals the product of all $\bar ax$, which is $\bar a^{\varphi(n)}$ times the same product. Cancel that unit.
\end{proof}

% t1-pass2 why
\paragraph{Why this works.}
Multiplying all the invertible residue classes by another invertible class merely rearranges them. So the product of all units is unchanged, even though the same operation multiplies that product by $a^{\varphi(n)}$. The product itself is invertible; cancelling it gives Euler's theorem. Before using this argument, ask why multiplication really is a permutation. The inverse of $a$ is exactly what makes the rearrangement reversible.

\subsection*{Division has an entrance requirement}
Modulo $7$, the class of $3$ has an inverse because $3\cdot5\equiv1$. Thus solving $3x\equiv2$ means multiplying by $5$, and the answer is $x\equiv3$. Try the same move modulo $6$: the congruence $3x\equiv3$ has several solutions, because $3$ has no inverse. Cancelling it would erase information. This example gives the rule to carry forward: compute $\gcd(c,m)$ before cancelling a factor $c$ modulo $m$.

The same issue controls exponents. Euler's theorem allows us to reduce exponents modulo $\varphi(n)$ \emph{only if the base is coprime to $n$}. For instance $\varphi(8)=4$, but $2^4\equiv0\pmod8$ while $2^0\equiv1\pmod8$. Reducing the exponent from $4$ to $0$ here would turn a false equality into an impressively short calculation.

\textbf{Repeated squaring.} Even without Euler's theorem, modular powers are quick. Since $13=8+4+1$, to compute $7^{13}\bmod23$ square repeatedly:
\[
\begin{aligned}
7^2&\equiv3,\qquad 7^4\equiv9,\qquad 7^8\equiv12,\\
7^{13}&\equiv12\cdot9\cdot7=756\equiv20\pmod{23}.
\end{aligned}
\]
At every step we reduce the result modulo $23$, so no enormous integer is needed. An exponent with $s$ binary digits requires at most $s-1$ squarings and $s-1$ further multiplications. This elementary idea is one of the workhorses of public-key cryptography.

\begin{think}
Compute $(p-1)!\bmod p$ for $p=5,7,11$. What do you notice? In $\Z/7\Z$, pair each of $1,\dots,6$ with its inverse: which numbers are left alone, and why can there be no others?
\end{think}

\begin{theorem}[Wilson]\label{thm:wilson}
For $p$ prime, $(p-1)!\equiv-1\pmod p$.
\end{theorem}
\begin{proof}
If $p=2$, the product is $1!=1\equiv-1\pmod2$, so assume $p$ is odd. In the field $\Z/p\Z$, $x^2=1$ means $(x-1)(x+1)=0$, so $x=\pm1$: only $1$ and $-1$ are their own inverses. The other non-zero classes pair off with their inverses, each pair having product $1$. So $(p-1)!\equiv1\cdot(-1)=-1$.
\end{proof}

% t1-pass2 why
\paragraph{Why this works.}
Pair every nonzero class modulo an odd prime with its inverse. Most classes form pairs whose product is $1$. Only the classes solving $x=x^{-1}$ stand alone. That equation is $x^2=1$, so in a field the two fixed points are $1$ and $-1$. Their product is $-1$. The other pairs contribute nothing new, giving Wilson's theorem. Checking the exceptional prime $2$ separately keeps the reasoning honest.

% t1-pass2 examples b31b4a00
\begin{examplex}[arithmetic modulo $12$]
One has $7+8=15\equiv 3$ and $7\cdot 8=56\equiv 8$ modulo $12$. The residue $5$ is a unit: $5\cdot 5=25\equiv 1$. The residue $4$ is not, because every multiple of $4$ is even and not $1$ modulo $12$. Figure~\ref{fig:clock} is this arithmetic drawn as hours.
\end{examplex}

\begin{examplex}[the last two digits of $7^4$]
Modulo $100$, $7^2=49$ and $7^4=49^2=2401\equiv 1$. So the order of $7$ modulo $100$ divides $4$ on this evidence, but Euler's theorem only promised exponent $\varphi(100)=40$. A direct power can be much smaller than the group order, which is why the order is defined as the least positive exponent in the next chapter.
\end{examplex}

\section{The Chinese remainder theorem}

% t1-pass2 motive 44b5e511
Suppose a number leaves remainder $2$ on division by $3$ and remainder $3$ on division by $5$. Try the numbers $2,5,8,11,\ldots$ that satisfy the first condition: $8$ satisfies both, and every solution differs from $8$ by a multiple of $15$. The Chinese remainder theorem turns this experiment into a formula. Coprime moduli carry independent information; B\'ezout's identity tells us how to fit their answers together. When they share a factor, we must first check compatibility.

\begin{think}
Draw a grid with $3$ rows and $4$ columns, and write each $x=0,1,\dots,11$ in row $x\bmod3$ and column $x\bmod4$. Does every cell receive exactly one number? Now try $2$ rows and $4$ columns with $x=0,\dots,7$. What property of the pair of moduli makes the difference?
\end{think}

A \emph{ring morphism} is a map preserving addition, multiplication and
$1$. A bijective ring morphism is an \emph{isomorphism}: it identifies two
presentations of the same arithmetic. Its inverse also preserves operations,
since the defining identities can be read backwards through a bijection.
The product of two rings uses coordinatewise operations.

\begin{theorem}[Chinese remainder theorem]\label{thm:crt}
If $m,n\ge1$ are coprime, the map $\Z/mn\Z\to\Z/m\Z\times\Z/n\Z$, $x\mapsto(x\bmod m,\ x\bmod n)$ is a ring isomorphism. Explicitly, if $mu+nv=1$, the solution of $x\equiv a\pmod m$, $x\equiv b\pmod n$ is $x\equiv anv+bmu\pmod{mn}$.
\end{theorem}
\begin{proof}
The map is a well-defined ring morphism. It is injective: if $m\mid x$ and $n\mid x$ with $m,n$ coprime then $mn\mid x$ (Gauss). Both sides have $mn$ elements, so it is bijective. For the formula: $nv\equiv1\pmod m$ and $nv\equiv0\pmod n$, and symmetrically for $mu$.
\end{proof}

Figure~\ref{fig:crt} shows the theorem as a walk: adding $1$ moves one step down and one step right, wrapping around the edges like on a torus. When the moduli are coprime, the walk visits every cell before coming back.

\begin{figure}[ht]
\centering
\fig{ch03-crt}
\caption{Left: $x\mapsto(x\bmod3,x\bmod4)$ for $x=0,\dots,11$. Each step $x\to x+1$ moves diagonally (the arrows show the steps that do not wrap around), and every cell is hit exactly once. Right: with the non-coprime moduli $2$ and $4$, the walk closes up after $\operatorname{lcm}(2,4)=4$ steps, so half the cells (grey) are never reached and the others are hit twice.}
\label{fig:crt}
\end{figure}

Units correspond to pairs of units, so $\varphi(mn)=\varphi(m)\varphi(n)$ for coprime $m,n$. With $\varphi(p^k)=p^k-p^{k-1}$ (the non-multiples of $p$), this recovers $\varphi(n)=n\prod_{p\mid n}\left(1-\frac1p\right)$ from Volume~1.

\begin{proposition}[linear congruences]\label{prop:lincong}
$ax\equiv b\pmod n$ has a solution if and only if $d=\gcd(a,n)$ divides $b$, and then it has exactly $d$ solutions modulo $n$.
\end{proposition}
\begin{proof}
Solvability means $b\in a\Z+n\Z=d\Z$. If $d\mid b$, divide through: $\frac adx\equiv\frac bd\pmod{\frac nd}$ has a unique solution $x_0$ modulo $\frac nd$, since $\frac ad$ is a unit there; it gives $d$ classes modulo $n$.
\end{proof}

\subsection*{Reassembling a number, and knowing when it is possible}
Solve $x\equiv2\pmod5$ and $x\equiv3\pmod7$. Since $3\cdot5-2\cdot7=1$, the theorem gives
\[
x\equiv2(-14)+3(15)=17\pmod{35}.
\]
Check both remainders: $17=3\cdot5+2=2\cdot7+3$. The two ingredients $-14$ and $15$ act as switches: one is $1$ modulo $5$ and $0$ modulo $7$, the other reverses those roles. This is why adding the two prescribed remainders works.

\textbf{When the moduli share a factor.} For arbitrary positive $m,n$, put $d=\gcd(m,n)$. The pair $x\equiv a\pmod m$, $x\equiv b\pmod n$ is solvable precisely when $d\mid b-a$, and then the answer is one class modulo $\operatorname{lcm}(m,n)=mn/d$. To prove this, write $x=a+mt$: the second condition becomes $mt\equiv b-a\pmod n$. The linear-congruence proposition gives exactly the compatibility condition and one class of $t$ modulo $n/d$, hence one class of $x$ modulo $mn/d$.

Thus $x\equiv1\pmod6$, $x\equiv3\pmod8$ is compatible and gives $x\equiv19\pmod{24}$, whereas replacing the second remainder by $2$ makes it impossible: an integer cannot be both odd and even. The Chinese remainder theorem is powerful; it does not negotiate with parity.

\subsection*{A counting identity for the next chapter}
The following identity will let us count elements of each multiplicative order.
It is a partition of a finite set, not an assumption about primitive roots.

\begin{proposition}[Gauss]\label{prop:phisum}
$\displaystyle\sum_{d\mid n}\varphi(d)=n$.
\end{proposition}
\begin{proof}
Write each of the $n$ fractions $\frac kn$, $1\le k\le n$, in lowest terms $\frac ad$: then $d\mid n$, $1\le a\le d$ and $\gcd(a,d)=1$, and each such pair arises exactly once. So $n=\sum_{d\mid n}\varphi(d)$.
\end{proof}

% t1-pass2 examples 44b5e511
\begin{examplex}[Sunzi's problem]
Solve $x\equiv 2\pmod 3$, $x\equiv 3\pmod 5$, $x\equiv 2\pmod 7$. First $x\equiv 2\pmod 3$ and $x\equiv 2\pmod 7$ give $x\equiv 2\pmod{21}$, because $3$ and $7$ are coprime and the two remainders agree. Then $x=21k+2\equiv 3\pmod 5$, so $k\equiv 1\pmod 5$ and $k=5t+1$. Thus $x=105t+23$. The smallest positive solution is $23$.
\end{examplex}

\begin{examplex}[when the moduli are not coprime]
The system $x\equiv 1\pmod 4$ and $x\equiv 2\pmod 6$ has no solution. A solution would give $1\equiv 2\pmod{\gcd(4,6)}$, that is $1\equiv 2\pmod 2$, which is false. The Chinese remainder theorem needs coprime moduli precisely so that this compatibility condition is automatic. When it holds, as for $x\equiv 1\pmod 4$ and $x\equiv 3\pmod 6$ (since $1\equiv 3\pmod 2$), solutions exist; $x=9$ is one.
\end{examplex}

\section{Lifting modulo prime powers}

% t1-pass2 motive 65285a19
A solution modulo $5$ is only a first approximation to a solution modulo $25$. For instance, $2^2+1=5$ is divisible by $5$ but not by $25$. Can we change $2$ by a multiple of $5$ to make the square plus $1$ divisible by $25$? Write the new candidate as $2+5t$, expand, and solve for $t$ modulo $5$. The derivative tells us why this linear correction has one solution when the root is simple. Repeating that step is Hensel lifting.

\begin{think}
Modulo $5$, $x=2$ solves $x^2\equiv-1$. Look for a solution modulo $25$ of the form $2+5t$. Expand the square: why does the term $25t^2$ not matter, and why is there exactly one good $t$ in $\{0,\dots,4\}$?
\end{think}

For $f(X)=\sum_{j=0}^d c_jX^j$, its \emph{formal derivative} is $f'(X)=\sum_{j=1}^d jc_jX^{j-1}$. No analytic limit is needed: the binomial theorem gives
\[
f(a+h)=f(a)+hf'(a)+h^2g(a,h)
\]
for a polynomial $g$ with integer coefficients. The constant and linear terms are displayed; every remaining term contains at least $h^2$. This elementary expansion is the entire engine of the next theorem.

\begin{theorem}[Hensel's lemma]\label{thm:hensel}
Let $f\in\Z[X]$, $p$ prime, $k\ge1$, and $a\in\Z$ with $f(a)\equiv0\pmod{p^k}$ and $f'(a)\not\equiv0\pmod p$. Then there is a unique $t\in\{0,\dots,p-1\}$ such that $f(a+tp^k)\equiv0\pmod{p^{k+1}}$.
\end{theorem}
\begin{proof}
Taylor's formula for polynomials gives $f(a+h)=f(a)+hf'(a)+h^2g(a,h)$ with $g$ having integer coefficients. With $h=tp^k$ and $2k\ge k+1$:
$f(a+tp^k)\equiv f(a)+tp^kf'(a)\pmod{p^{k+1}}$. Writing $f(a)=cp^k$, the condition becomes $c+tf'(a)\equiv0\pmod p$, which has exactly one solution $t$ modulo $p$ since $f'(a)$ is invertible modulo $p$.
\end{proof}

% t1-pass2 why
\paragraph{Why this works.}
When you replace a proposed root $a$ by $a+tp^k$, expand the polynomial. Modulo $p^{k+1}$, every term containing two or more factors $tp^k$ disappears; only the constant value and the linear correction survive. The linear coefficient is $f'(a)$. If it is invertible modulo $p$, exactly one value of the new digit $t$ works. That is why a simple root lifts uniquely. A multiple root may split or die instead.

Iterating, a simple root modulo $p$ lifts to a unique root modulo every $p^k$. This is Newton's method in the $p$-adic world (Volume~13).

\subsection*{A root has a family tree}
The equation $x^2+1\equiv0\pmod5$ has roots $2$ and $3$. To lift the root $2$ to a solution modulo $25$, try $x=2+5t$. Then $(2+5t)^2+1\equiv5+20t\pmod{25}$. Divide by $5$ and solve $1+4t\equiv0\pmod5$, giving $t\equiv1$ and the lifted root $7$. The correction works uniquely because the derivative $2x$ is not divisible by $5$ at the original root. A picture of branches is useful here only after this arithmetic is clear.

\begin{figure}[ht]
\centering\fig{ch03-hensel-branches}
\caption{Two simple roots of $X^2+1$ lifting through powers of $5$. The integers at each level are representatives, not increasing real approximations: $57$ and $68$ both square to $-1$ modulo $125$. Each arrow records reduction modulo the preceding power of $5$.}
\label{fig:hensel-branches}
\end{figure}

If the derivative vanishes modulo $p$, uniqueness can fail in two opposite ways. For $X^2$ the root $0$ modulo $5$ has \emph{five} lifts modulo $25$, namely $0,5,10,15,20$. For $X^2-5$, the same root modulo $5$ has \emph{no} lift modulo $25$, since $(5t)^2-5\equiv-5\not\equiv0\pmod{25}$. So the derivative hypothesis is not decoration: remove it and a root may multiply or disappear.

\subsection*{A compact strategy for congruence problems}
First decide whether the task concerns an additive equation, a product, or an exponent. For an equation $ax\equiv b$, compute $\gcd(a,n)$ before dividing. For a composite modulus, split into coprime prime powers and reassemble by the Chinese remainder theorem. For exponents, check coprimality before using Euler, and use repeated squaring to perform the remaining computation. For polynomial roots modulo $p^k$, inspect the derivative before promising a unique lift. These are different tools with different hypotheses; choosing the tool is often most of the mathematics.

% t1-pass2 examples 65285a19
\begin{examplex}[lifting the root of $X^2+1$ modulo $5$]
Modulo $5$ the roots are $2$ and $3$. From $2$, the next digit satisfies $2+5t$ with $(2+5t)^2\equiv -1\pmod{25}$, so $20t\equiv -5\pmod{25}$ and $t\equiv 1\pmod 5$. The lift is $7$, and $7^2=49\equiv -1\pmod{25}$. The other branch starts at $3$ and is $-7\equiv 18\pmod{25}$.
\end{examplex}

\begin{examplex}[a singular root that dies]
Modulo $5$, $X^2-5$ vanishes at $0$. Modulo $25$ the same congruence asks for a square equal to $5$, and no square is $5$ modulo $25$. The derivative $2X$ also vanishes at $0$, so the unique-lift machine refuses to start. Both behaviours, a unique child and no child, are the reason the derivative is a hypothesis rather than a convenience.
\end{examplex}

\begin{beauty}
Where this leads. The Chinese remainder theorem is the exact sieve identity of Chapter~\ref{ch:characters} once it is written with the M\"obius function, and it is the gluing of local solutions in Volume~12. Hensel's lemma, kept here for congruences, is the construction of $\Z_p$ in Volume~13. Euler's theorem is the reason a character sum over a complete set of units can be evaluated in Volume~5.
\end{beauty}

\subsection*{Proof workshop: what cancellation actually requires}

Modulo $12$, solve $5x\equiv7$. Since
$5\cdot5=25\equiv1\pmod{12}$, multiply both sides by $5$:
\[
x\equiv5\cdot7=35\equiv11\pmod{12}.
\]
Check by substitution: $5\cdot11=55\equiv7$. This is one
solution class, not just the chosen representative $11$.

Now change only the coefficient: $6x\equiv6\pmod{12}$.
The solutions include $x\equiv1$ and $x\equiv3$, because
$6$ times either is $6$ modulo $12$. Cancelling $6$ and
claiming $x\equiv1$ would discard a valid solution. The
general rule is not that cancellation sometimes feels unsafe:
you may cancel a factor modulo $m$ precisely when that factor
is a unit, equivalently when it is coprime to $m$.

\begin{figure}[ht]
\centering
\fig{ch03-unit-classes}
\caption{Every residue modulo $12$ has a position, but only
$1,5,7,11$ have multiplicative inverses. A cancellation step
requires a unit; the other classes cannot be divided out.}
\label{fig:congruence-unit-classes}
\end{figure}


\section{Exercises}

\begin{exercise}[$\star$]\label{ex:last2}
Find the last two digits of $7^{2026}$.
\end{exercise}

\begin{exercise}[$\star$]\label{ex:sunzi}
\emph{Sunzi's remainder problem.} Find the integers that leave remainders $2$, $3$, $2$ when divided by $3$, $5$, $7$.
\end{exercise}

\begin{exercise}[$\star$]\label{ex:digits}
Show that a nonnegative integer is congruent modulo $9$ to the sum of its digits, and modulo $11$ to the alternating sum of its digits (starting from the units). State and prove the corresponding rule for negative integers.
\end{exercise}

\begin{exercise}[$\star\star$]\label{ex:sqrt1}
\begin{enumerate}[label=(\alph*)]
\item Let $p$ be an odd prime and $k\ge1$. Show that the solutions of $x^2\equiv1\pmod{p^k}$ are $x\equiv\pm1$.
\item Let $n$ be odd with $r$ distinct prime factors. How many solutions does $x^2\equiv1\pmod n$ have?
\end{enumerate}
\end{exercise}

\begin{exercise}[$\star\star$]\label{ex:wilsonconv}
Show that if $n>4$ is composite then $(n-1)!\equiv0\pmod n$. Deduce that $n>1$ is prime if and only if $(n-1)!\equiv-1\pmod n$.
\end{exercise}

\begin{exercise}[$\star\star$]\label{ex:hensel}
Find $x\in\{0,\dots,124\}$ with $x^2\equiv-1\pmod{125}$, by lifting a solution modulo $5$.
\end{exercise}

\begin{exercise}[$\star\star$]\label{ex:tower}
Show that the last two digits of $3^{3^{3^3}}$ are $87$.
\end{exercise}

\begin{exercise}[$\star\star$]\label{ex:korselt}
\emph{Korselt's criterion, one direction.} Let $n$ be squarefree and composite, and suppose $p-1\mid n-1$ for every prime $p\mid n$. Show that $a^n\equiv a\pmod n$ for every integer $a$. Check that $561=3\cdot11\cdot17$ qualifies: it is a Carmichael number (Volume~1, the sums and induction chapter).
\end{exercise}

\begin{exercise}[$\star\star\star$]\label{ex:lucas}
\emph{Lucas' theorem.} Let $p$ be prime, and write $n=\sum_in_ip^i$ and $k=\sum_ik_ip^i$ in base $p$. Show that
\[
\binom nk\equiv\prod_i\binom{n_i}{k_i}\pmod p,
\]
with $\binom{n_i}{k_i}=0$ when $k_i>n_i$. \emph{Hint: in $(\Z/p\Z)[X]$, $(1+X)^p=1+X^p$.} Deduce that $\binom nk$ is odd if and only if each binary digit of $k$ is at most the corresponding digit of $n$, which explains the Sierpi\'nski triangle of Volume~1.
\end{exercise}

\begin{exercise}[$\star\star$]\label{ex:t1-toy-rsa}
A miniature RSA system uses $N=55=5\cdot11$, public exponent $e=3$ and private exponent $d=27$. Show that encrypting a residue $m$ by $c\equiv m^3\pmod{55}$ and then decrypting by $c^{27}$ returns $m$ for \emph{every} residue $m$, including non-units. Encrypt $m=12$ and decrypt the result. Explain why Euler's theorem alone, applied modulo $55$, does not prove the assertion for all messages.
\end{exercise}

\begin{exercise}[$\star\star\star$]\label{ex:t1-two-power-roots}
Find the solutions of $x^2\equiv1\pmod{2^k}$ for $k=1,2$. Prove that for every $k\ge3$ there are exactly four classes,
\[
1,\quad -1,\quad 1+2^{k-1},\quad -1+2^{k-1}\pmod{2^k}.
\]
Where does the argument for odd primes in Exercise~\ref{ex:sqrt1} break down?
\end{exercise}

\begin{exercise}[$\star$]\label{ex:t1-sunzi23}
Find the smallest positive integer $x$ with $x\equiv 2\pmod 3$, $x\equiv 3\pmod 5$ and $x\equiv 2\pmod 7$.
\end{exercise}

\begin{beauty}
The Chinese remainder theorem says that arithmetic modulo $n$ is arithmetic modulo each prime power of $n$, independently. It speeds up RSA decryption by working with smaller moduli (the gain is roughly a factor of four under a schoolbook arithmetic cost model), and in Volume~2 it becomes a statement about products of rings. Hensel's lemma is the door to the $p$-adic numbers of Volume~13, where $\sqrt{-1}$ exists as a genuine $5$-adic number, the limit of $2,7,57,\dots$ from Exercise~\ref{ex:hensel}.
\end{beauty}

\begin{exercise}[$\star\star$]\label{pb:primepoly}
Let $P\in\Z[X]$ be non-constant.
\begin{enumerate}[label=(\alph*)]
\item Show that there are infinitely many integers $n\ge0$ for which $|P(n)|$ is not a prime number.
\item Show that infinitely many primes divide at least one of the numbers $P(0),P(1),P(2),\dots$
\end{enumerate}
\end{exercise}

\begin{exercise}[$\star\star$]\label{pb:fpfunctions}
Let $p$ be a prime.
\begin{enumerate}[label=(\alph*)]
\item Show that every map $f:\mathbb F_p\to\mathbb F_p$ is the map $x\mapsto P(x)$ of exactly one polynomial $P\in\mathbb F_p[X]$
of degree at most $p-1$.
\item Show that the coefficient of $X^{p-1}$ in this polynomial is $-\sum_{a\in\mathbb F_p}f(a)$.
\item Suppose $p$ odd and $f$ a bijection. Show that $P$ has degree at most $p-2$, and that for each $t$ with $1\le t\le p-2$
the polynomial of the map $x\mapsto f(x)^t$ has degree at most $p-2$ too. Check this for $f(x)=x^3$ on $\mathbb F_5$.
\end{enumerate}
\end{exercise}

\begin{exercise}[$\star\star$]\label{pb:aks}
Let $n\ge2$ and let $a$ be an integer prime to $n$. Show that $n$ is prime if and only if
\[
(X+a)^n\equiv X^n+a\pmod n
\]
in $\Z[X]$, that is, coefficient by coefficient.
\end{exercise}

\begin{exercise}[$\star\star$]\label{pb:wolsten}
\emph{Wolstenholme's theorem.} Let $p\ge5$ be prime.
\begin{enumerate}[label=(\alph*)]
\item Show that $\sum_{k=1}^{p-1}\frac1{k^2}\equiv0\pmod p$.
\item Show that $1+\frac12+\dots+\frac1{p-1}\equiv0\pmod{p^2}$.
\item Show that $\binom{2p-1}{p-1}\equiv1\pmod{p^3}$, that is, $\binom{2p}p\equiv2\pmod{p^3}$.
\end{enumerate}
\end{exercise}

\begin{exercise}[$\star\star$]\label{pb:jacobsthal}
Let $p$ be prime and $a\ge b\ge0$. Show that $\binom{ap}{bp}\equiv\binom ab\pmod{p^2}$.
\end{exercise}

\begin{exercise}[$\star\star$]\label{pb:morley}
Let $p$ be an odd prime. Show that
\[
\binom{p-1}{(p-1)/2}\equiv(-1)^{(p-1)/2}\,4^{p-1}\pmod{p^2}.
\]
\end{exercise}

\begin{exercise}[$\star\star\star$]\label{pb:glaisher}
Let $p$ be an odd prime and $B_{p-1}$ the Bernoulli number. Show that $pB_{p-1}$ has no $p$ in its denominator and that
\[
(p-1)!\equiv pB_{p-1}-p\pmod{p^2}.
\]
\end{exercise}

\begin{exercise}[Journey]\label{pb:ramanujan5}
\emph{Ramanujan's congruence.} We prove that $p(5n+4)\equiv0\pmod5$ for every $n\ge0$. Fix $q$ with $0<|q|<1$ and let, for
$z\in\C^*$,
\[
H(z)=\prod_{n\ge1}(1+zq^{2n-1})(1+z^{-1}q^{2n-1})=\sum_{m\in\Z}c_mz^m
\]
(a Laurent series, the coefficients depending on $q$).
\begin{enumerate}[label=(\alph*)]
\item Show that $H(q^2z)=\frac{H(z)}{qz}$, and deduce that $c_m=q^{m^2}c_0$ for every $m\in\Z$.
\item By evaluating at $z=i$, and at $z=-1$ with $q$ replaced by $q^4$, show that $c_0(q)E(q^2)=c_0(q^4)E(q^8)$, and deduce the
\emph{Jacobi triple product}
\[
\prod_{n\ge1}(1-q^{2n})(1+zq^{2n-1})(1+z^{-1}q^{2n-1})=\sum_{m\in\Z}q^{m^2}z^m.
\]
\item Deduce \emph{Jacobi's identity} $E(x)^3=\sum_{k\ge0}(-1)^k(2k+1)x^{k(k+1)/2}$ for $|x|<1$.
\item Show that in $xE(x)^4=xE(x)\cdot E(x)^3$, every coefficient of a power $x^{5m}$ is divisible by $5$.
\item Using $E(x)^5\equiv E(x^5)\pmod5$, conclude.
\end{enumerate}
\end{exercise}

\begin{exercise}[$\star\star\star$]\label{pb:kloosterman}
\emph{Kloosterman sums.} Let $p$ be an odd prime. For $a,b\in\mathbb F_p$ let
\[
K(a,b)=\sum_{x\in\mathbb F_p^\times}e\Bigl(\frac{ax+b\bar x}p\Bigr),
\]
where $\bar x$ is the inverse of $x$ modulo $p$, and $K(a)=K(a,1)$.
\begin{enumerate}[label=(\alph*)]
\item Show that $K(a)$ is real, that $K(0)=-1$, and that $K(a,b)=K(ab)$ when $b\ne0$.
\item Show that $\sum_{a\in\mathbb F_p}K(a)^2=p^2-p$.
\item Show that $\sum_{a\in\mathbb F_p}K(a)^4=2p^3-3p^2-3p$.
\item Deduce that $|K(a)|<1.19\,p^{3/4}$ for every $a$. How does this compare with the trivial bound, and with what (b) gives?
\end{enumerate}
\end{exercise}

\section{Hints}

\begin{hint}{ex:last2}
Think of the last two digits as a residue modulo $100$. Before applying Euler's theorem, check the smaller powers $7^2$ and $7^4$. A short period makes the exponent $2026$ manageable without ever constructing the large number.
\end{hint}
\begin{hint}{ex:sunzi}
Solve the conditions one at a time. The numbers congruent to $3$ modulo $5$ have a simple pattern; which one also leaves remainder $2$ modulo $3$? Now vary it by a multiple of $15$ until the condition modulo $7$ holds. Check your final answer in all three moduli.
\end{hint}
\begin{hint}{ex:digits}
Write the integer as $\sum_i d_i10^i$. Modulo $9$, the base $10$ acts like $1$; modulo $11$, it acts like $-1$. Substitute those two values before simplifying. The digit rules follow from the polynomial expression, not from coincidence.
\end{hint}
\begin{hint}{ex:sqrt1}
Factor $x^2-1$ as $(x-1)(x+1)$. For an odd prime $p$, why can $p$ not divide both factors? Once you know where the whole prime power must go, solve independently modulo each prime power and use the Chinese remainder theorem to combine the sign choices.
\end{hint}
\begin{hint}{ex:wilsonconv}
For a composite $n=ab$, the factorial $(n-1)!$ often contains both $a$ and $b$ as separate factors. The only awkward case is $a=b$: look for the distinct factors $a$ and $2a$. Verify the inequality needed for $2a<n$, and check the smallest exceptional square separately.
\end{hint}
\begin{hint}{ex:hensel}
Start from $x\equiv2\pmod5$ and apply the lifting step of Theorem~\ref{thm:hensel} twice, exactly as in the last thinking prompt.
\end{hint}
\begin{hint}{ex:tower}
Modulo $100$, a power of $3$ only depends on its exponent modulo $\varphi(100)=40$. So first find $3^{27}\bmod40$; note that $3^4=81$.
\end{hint}
\begin{hint}{ex:korselt}
By the Chinese remainder theorem it is enough to prove $a^n\equiv a\pmod p$ for each prime $p\mid n$. Treat $p\mid a$ and $p\nmid a$ separately, and use Fermat with $n-1=(p-1)q$.
\end{hint}
\begin{hint}{ex:lucas}
In $(\Z/p\Z)[X]$, write
\[(1+X)^n=\prod_i\left((1+X)^{p^i}\right)^{n_i}=\prod_i\left(1+X^{p^i}\right)^{n_i},\]
and find the coefficient of $X^k$ on both sides. Uniqueness of base-$p$ expansions does the rest.
\end{hint}

\begin{hint}{ex:t1-toy-rsa}
The composed exponent is $81$. Work separately modulo $5$ and modulo $11$, distinguishing a zero residue from a non-zero one. For the numerical decryption, reduce the exponent modulo $20$ or use repeated squaring.
\end{hint}
\begin{hint}{ex:t1-two-power-roots}
A root is odd. Of the two consecutive even integers $x-1,x+1$, one is divisible by $4$ and the other has exactly one factor $2$. Use this to force $x\equiv1$ or $-1\pmod{2^{k-1}}$, then list the two lifts of each class modulo $2^k$.
\end{hint}

\begin{hint}{ex:t1-sunzi23}
The first and third congruences together say $x\equiv 2\pmod{21}$. Then solve one linear congruence modulo $5$.
\end{hint}

\begin{hint}{pb:primepoly}
(a) If $P(a)=\pm p$, look at $P(a+kp)$ modulo $p$. (b) Evaluate $P$ at a multiple of $P(0)$ times the product of the primes
you already have.
\end{hint}

\begin{hint}{pb:fpfunctions}
$1-(x-a)^{p-1}$ is the indicator of $\{a\}$ on $\mathbb F_p$. For (c), use the power sums $\sum_xx^t$ modulo $p$.
\end{hint}

\begin{hint}{pb:aks}
If $q^k$ is the exact power of a prime $q$ dividing $n$, look at the coefficient of $X^q$.
\end{hint}

\begin{hint}{pb:wolsten}
(a) $k\mapsto\frac1k$ permutes the non-zero residues. (b) Pair $k$ with $p-k$. (c) $\binom{2p-1}{p-1}=\prod_{k=1}^{p-1}\bigl(1+\frac pk\bigr)$.
\end{hint}

\begin{hint}{pb:jacobsthal}
Vandermonde with $ap=(a-1)p+p$, and induction on $a$.
\end{hint}

\begin{hint}{pb:morley}
$\binom{p-1}k=\prod_{j\le k}\bigl(\frac pj-1\bigr)$, and relate $H_{(p-1)/2}$ to $q_p(2)$ through Problem~\ref{pb:fermatquot}(b).
\end{hint}

\begin{hint}{pb:glaisher}
Faulhaber's formula for $\sum_{k<p}k^{p-1}$ modulo $p^2$, von Staudt--Clausen, and Problem~\ref{pb:fermatquot}(c).
\end{hint}

\begin{hint}{pb:ramanujan5}
(b) At $z=i$ the odd powers cancel. (c) Put $q=x^{1/2}$, $z=-wx^{1/2}$, divide by $1-w^{-1}$ and let $w\to1$. (d) Multiply the
exponent by $8$ and complete the squares modulo $5$.
\end{hint}

\begin{hint}{pb:kloosterman}
(a) Substitute $x\mapsto-x$, and $x\mapsto bx$. (b) Sum over $a$ first. (c) Do not attack $K(a)^4$ directly: compute $\sum_{a,b}|K(a,b)|^4$ in two
ways, once with (a) and once by summing over $a$ and $b$ first, which counts solutions of $x_1+x_2=x_3+x_4$, $\bar x_1+\bar x_2=\bar x_3+\bar x_4$.
(d) One term of a sum of non-negative terms is at most the sum.
\end{hint}

\section{Solutions}

\begin{solution}{ex:last2}
\idea The last two digits of an integer are its residue modulo $100$. Do not calculate a towering exponent. First look for a short period of the base: $7^2\equiv49$ and $7^4\equiv1\pmod{100}$. Once a power returns to $1$, only the exponent modulo that period matters. Check the base is invertible before using a period theorem; here it is, since $\gcd(7,100)=1$.
\writeup We need only a remainder modulo $100$, not the decimal expansion of $7^{2026}$. Compute $7^2=49$ and $7^4=2401\equiv1\pmod{100}$. Since $2026=4\cdot506+2$, we can group the exponent:
\[
7^{2026}=(7^4)^{506}7^2\equiv49\pmod{100}.
\]
Thus the final two digits are $49$. The small return exponent is the reason this computation stays easy.
\end{solution}

\begin{solution}{ex:sunzi}
\idea The moduli $3$, $5$ and $7$ are pairwise coprime, so there will be exactly one solution class modulo their product $105$. For small moduli, solve two conditions first by listing a few candidates, then impose the third condition on that arithmetic progression. At the end, verify the answer in all three original congruences; the Chinese remainder theorem promises uniqueness, not immunity from arithmetic errors.
\writeup First solve the conditions modulo $3$ and $5$. Among the numbers $3,8,13$ congruent to $3$ modulo $5$, only $8$ is $2$ modulo $3$, so $x\equiv8\pmod{15}$. Write $x=8+15t$. Modulo $7$, this is $x\equiv1+t$, and the last condition gives $t\equiv1\pmod7$. Therefore
\[
x\equiv23\pmod{105},\qquad x=23+105k,\quad k\in\Z.
\]
The Chinese remainder theorem guarantees that this is the complete family because the three moduli are pairwise coprime.
\end{solution}

\begin{solution}{ex:digits}
\idea Write a number in base ten as a polynomial in $10$, with its digits as coefficients. Modulo $9$, the base $10$ becomes $1$, so every place value becomes $1$. Modulo $11$, it becomes $-1$, so the signs alternate. The familiar tests for divisibility by $9$ and $11$ are not tricks: they are evaluations of the same polynomial at different residues.
\writeup Write a nonnegative integer as $n=\sum_i d_i10^i$, with $d_i$ its decimal digits. Modulo $9$, the base is $10\equiv1$, so $n\equiv\sum_i d_i$. Modulo $11$, the base is $10\equiv-1$, so $n\equiv\sum_i(-1)^id_i$. These are the digit-sum and alternating-digit-sum rules, obtained by legitimate substitution in congruences.

For a negative integer $-n$, apply the formulas to $n=|-n|$ and negate the result. For example, $-12\equiv-3\pmod9$, not $3\pmod9$. The sign of the integer must not be silently discarded.
\end{solution}

\begin{solution}{ex:sqrt1}
\idea Factor before counting roots: $x^2-1=(x-1)(x+1)$. Modulo an odd prime power, the two factors cannot both be divisible by the prime, since their difference is $2$. Thus all the required prime power must lie in one factor. For a modulus with several coprime prime-power factors, solve the equation independently in each component and use the Chinese remainder theorem to assemble the roots.
\writeup (a) If $x^2\equiv1\pmod{p^k}$, then $p^k\mid(x-1)(x+1)$. The odd prime $p$ cannot divide both factors because their difference is $2$. One factor is therefore coprime to $p^k$, so Gauss's lemma forces $p^k$ to divide the other. We obtain the two distinct classes $x\equiv\pm1\pmod{p^k}$.

(b) Write the odd modulus as $n=\prod_{i=1}^r p_i^{a_i}$. A root modulo $n$ is equivalent, by the Chinese remainder theorem, to an independent choice of a root modulo each prime power. Part (a) gives two choices at every factor. Every sign pattern is possible, so there are exactly $2^r$ roots.
\end{solution}

\begin{solution}{ex:wilsonconv}
\idea Wilson's theorem has a sharp converse, but the composite cases need care. If $n=ab$ has distinct proper factors $a,b$, both occur in $(n-1)!$, so their product divides it. When $n=a^2$, the two factors coincide and cannot be taken twice from the factorial; look instead for $a$ and $2a$. This works when $a>2$. The small exceptional square must be checked separately.
\writeup Suppose $n>4$ is composite and write $n=ab$ with $1<a\le b<n$. If $a<b$, the two distinct integers $a,b$ occur in $(n-1)!$, hence their product $n$ divides the factorial. If $a=b$, then $n=a^2$ with $a>2$, and the distinct integers $a$ and $2a$ both occur in $(n-1)!$. Their product $2n$ again makes the factorial divisible by $n$.

Thus every composite $n>4$ gives $(n-1)!\equiv0\pmod n$. For the exceptional composite $4$, $3!=6\equiv2\pmod4$, also not $-1$. Wilson's theorem gives the converse for primes. We have proved the claimed equivalence for all $n>1$.
\end{solution}

\begin{solution}{ex:hensel}
\idea Hensel's lemma is constructive: at each step, write the candidate $a+tp^k$, expand modulo $p^{k+1}$, and solve the resulting linear equation for $t$ modulo $p$. Start from $2^2\equiv-1\pmod5$.
\writeup Start with $f(X)=X^2+1$. Modulo $5$, $2$ is a root and $f'(2)=4$ is invertible. To lift it modulo $25$, try $x=2+5t$:
\[
f(x)\equiv5+20t\pmod{25}.
\]
The required equation is $1+4t\equiv0\pmod5$, so $t\equiv1$ and $x=7$.

For the next lift write $x=7+25t$. Modulo $125$,
$f(x)\equiv50+350t$. Dividing by $25$ yields $2+14t\equiv0\pmod5$, hence $t\equiv2$ and $x=57$. Direct verification gives $57^2+1=3250=26\cdot125$. This is the unique lift of the root $2$ modulo $5$; the other initial root $3$ lifts separately.
\end{solution}

\begin{solution}{ex:tower}
\idea Work from the top of the tower down: modulo $100$, a power of $3$ depends on its exponent modulo $40=\varphi(100)$ (Euler); and the exponent $3^{27}$ modulo $40$ depends on $27$ modulo the order of $3$ modulo $40$, which is $4$. Each level reduces the modulus.
\writeup Put $E=3^{27}$. We want $3^E$ modulo $100$. Since $\gcd(3,100)=1$ and $\varphi(100)=40$, Euler's theorem says it is enough to know $E$ modulo $40$. But $3^4=81\equiv1\pmod{40}$, so $E=3^{27}\equiv3^3=27\pmod{40}$. Hence $3^E\equiv3^{27}\pmod{100}$. Finally $3^{10}\equiv49$ and $3^{20}\equiv49^2=2401\equiv1\pmod{100}$, giving $3^{27}\equiv3^7=2187\equiv87\pmod{100}$. The tower's last two digits are $87$.
\end{solution}

\begin{solution}{ex:korselt}
\idea To prove a congruence modulo a squarefree $n$, prove it modulo each prime factor $p$ (Chinese remainder theorem). Modulo $p$, Fermat says $a^{p-1}\equiv1$, so $a^{n-1}\equiv1$ as soon as $p-1$ divides $n-1$; that is exactly the hypothesis.
\writeup Fix a prime factor $p$ of $n$. The hypothesis gives $p-1\mid n-1$. If $p\nmid a$, Fermat's little theorem implies $a^{n-1}\equiv1\pmod p$, so $a^n\equiv a\pmod p$. If $p\mid a$, both sides are $0$ modulo $p$, and the same conclusion holds. Thus every prime factor of $n$ divides $a^n-a$. Since $n$ is squarefree, their product $n$ divides $a^n-a$ as well. For $561=3\cdot11\cdot17$, the number $560$ is divisible by $2$, $10$ and $16$, so the criterion applies.
\end{solution}

\begin{solution}{ex:lucas}
\idea Binomial coefficients are coefficients of $(1+X)^n$, so compute $(1+X)^n$ modulo $p$ cleverly. The ``freshman's dream'' $(1+X)^p=1+X^p$ in characteristic $p$ splits $(1+X)^n$ along the base-$p$ digits of $n$; then reading the coefficient of $X^k$ amounts to writing $k$ in base $p$.
\writeup Work in the polynomial ring over $\mathbb F_p$. Since the intermediate coefficients $\binom pj$ are divisible by $p$, we have $(1+X)^p=1+X^p$, and by iteration $(1+X)^{p^i}=1+X^{p^i}$. If $n=\sum_i n_ip^i$ then
\[
(1+X)^n=\prod_i(1+X^{p^i})^{n_i}
=\prod_i\sum_{j=0}^{n_i}\binom{n_i}{j}X^{jp^i}.
\]
To obtain exponent $k=\sum_i k_ip^i$, uniqueness of base-$p$ notation forces $j=k_i$ in each factor. The coefficient is $\prod_i\binom{n_i}{k_i}$, with value zero when some $k_i>n_i$. Comparing with the coefficient $\binom nk$ on the left proves Lucas's theorem. In base $2$, this says the binomial coefficient is odd precisely when every $1$ digit of $k$ also occurs in $n$.
\end{solution}

\begin{solution}{ex:t1-toy-rsa}
\idea Euler works for units, so split the modulus until every non-zero residue is a unit in a field. Fermat handles those residues, and zero handles itself; the Chinese remainder theorem reassembles the answer.
\writeup Encryption and decryption together raise the message $m$ to exponent $3\cdot27=81$. Modulo $5$, Fermat's theorem gives $m^4\equiv1$ for nonzero $m$ and $81=20\cdot4+1$; for $m=0$, the desired congruence is automatic. Thus $m^{81}\equiv m\pmod5$ for every residue. The same proof modulo $11$ uses exponent $10$ and $81=8\cdot10+1$. The Chinese remainder theorem combines the two claims modulo $55$.

For $m=12$, encryption gives $12^3=1728\equiv23\pmod{55}$. Decryption uses $23^2\equiv34$ and $23^4\equiv1$, so $23^{27}\equiv23^3\equiv34\cdot23=782\equiv12\pmod{55}$. The toy parameters show correctness, not secrecy.
\end{solution}

\begin{solution}{ex:t1-two-power-roots}
\idea The factors $x-1,x+1$ differ by $2$, so they can share a factor $2$, but no higher power. Track which factor owns the remaining powers of $2$.
\writeup For $k=1$ the only root is $1$; for $k=2$ the roots are $1$ and $3$. Suppose $k\ge3$ and $2^k\mid(x-1)(x+1)$. Then $x$ is odd. The two even factors differ by $2$, so exactly one contains just one factor of $2$, while the other must be divisible by $2^{k-1}$. Thus $x\equiv\pm1\pmod{2^{k-1}}$.

Each parent class has two lifts modulo $2^k$, namely $x=\varepsilon+t2^{k-1}$ with $\varepsilon=\pm1$ and $t=0,1$. Direct expansion gives
\[
x^2-1=\varepsilon t2^k+t^22^{2k-2},
\]
divisible by $2^k$ when $k\ge2$. These four classes are distinct for $k\ge3$, so they are all the roots. The exceptional behaviour comes from the prime $2$ dividing both $x-1$ and $x+1$.
\end{solution}

\begin{solution}{ex:t1-sunzi23}
\idea Coprime moduli may be combined one pair at a time.
\writeup The two conditions modulo $3$ and $7$ both prescribe remainder $2$, so together they give $x\equiv2\pmod{21}$. Write $x=21k+2$. The remaining condition modulo $5$ becomes $k+2\equiv3$, hence $k\equiv1\pmod5$. Therefore
\[
x=105t+23,\qquad t\in\Z.
\]
The smallest positive solution is $23$. It has remainders $2$, $3$ and $2$ on division by $3$, $5$ and $7$, respectively.
\end{solution}

\begin{solution}{pb:primepoly}
\idea A polynomial with integer coefficients respects congruences: $P(a+kp)\equiv P(a)\pmod p$. So once $P(a)=\pm p$, all the
values $P(a+kp)$ are multiples of $p$; if they were all prime they would all be $\pm p$, too many times for a non-constant
polynomial. For (b), Euclid's trick: build a value of $P$ that is $\equiv P(0)$ modulo a big product, then divide out
$P(0)$.

\writeup (a) Suppose a nonconstant polynomial $P\in\Z[X]$ took prime absolute values at every sufficiently large nonnegative integer. Choose one such input $a$ and set $p=|P(a)|$. Every integer $a+kp$ has the same residue as $a$ modulo $p$, so
\[
P(a+kp)\equiv P(a)\equiv0\pmod p.
\]
The absolute value is assumed prime, hence must equal $p$. For infinitely many $k$, the polynomial value therefore belongs to the two-element set $\{p,-p\}$. One of the nonzero polynomials $P-p$ and $P+p$ would have infinitely many roots. That is impossible, so $P$ cannot take only prime absolute values eventually.

(b) Write $P(X)=c+XR(X)$, where $c=P(0)$. If $c=0$, every prime $p$ divides $P(p)$ and there is nothing more to prove. Otherwise assume the only primes dividing nonzero values of $P$ are $p_1,\ldots,p_r$. Put $m=cp_1\cdots p_r$. Then
\[
P(mt)=c\bigl(1+p_1\cdots p_r\,tR(mt)\bigr).
\]
The expression in parentheses is a nonconstant integer polynomial in $t$, and so has an absolute value greater than $1$ for some integer $t$. Let $q$ be one of its prime factors. It also divides $P(mt)$, so by assumption $q=p_i$ for some $i$. But the parenthesised factor is congruent to $1$ modulo every $p_i$. This contradiction proves that infinitely many primes divide values of $P$.

\tools Volume~2 (primes and congruences); Volume~1 (nonzero polynomials have finitely many roots).
\leads This is Schur's theorem. The distribution of primes dividing polynomial values is a later topic; whether $n^2+1$ itself is prime for infinitely many $n$ remains open.
\end{solution}

\begin{solution}{pb:fpfunctions}
\idea The exponent $p-1$ is special in $\mathbb F_p$: by Fermat, $y^{p-1}$ is $1$ for $y\ne0$ and $0$ for $y=0$. So
$1-(x-a)^{p-1}$ is the indicator of the point $a$, and any function is a combination of indicators: this is Lagrange
interpolation in its most economical form. Its top coefficient is read off at once, and a bijection has the same multiset of
values as the identity.

\writeup (a) For each $a\in\mathbb F_p$, define the polynomial
\[
\delta_a(X)=1-(X-a)^{p-1}.
\]
Fermat's little theorem says that $\delta_a(x)$ equals $1$ at $x=a$ and $0$ at every other field element. Therefore
\[
P(X)=\sum_{a\in\mathbb F_p}f(a)\delta_a(X)
\]
agrees with the given function $f$ everywhere and has degree at most $p-1$. If another polynomial of degree at most $p-1$ agreed with $f$, their difference would have $p$ distinct roots and degree at most $p-1$, so it would be zero. Existence and uniqueness are both proved.

(b) The leading coefficient of $\delta_a$ at degree $p-1$ is $-1$. Consequently that of $P$ is $-\sum_af(a)$.

(c) If $f$ is a bijection and $p$ is odd, its values run through the field, so $\sum_af(a)=\sum_{x\in\mathbb F_p}x=0$. Part (b) then gives $\deg P\le p-2$. More generally, for $1\le t\le p-2$, the sum $\sum_af(a)^t$ equals the power sum $\sum_xx^t$, which vanishes modulo $p$. Applying (b) to the function $f^t$ again bounds its representing polynomial degree by $p-2$.

For $f(x)=x^3$ over $\mathbb F_5$, the map is a bijection since $3$ is coprime to $4$. Reducing powers with $x^5=x$, we obtain representative polynomials $X^3$, $X^2$ and $X$ for $f,f^2,f^3$; their degrees are at most $3$.

\tools Volume~2 (Fermat and power sums); Volume~1 (polynomial root bound).
\leads These are the elementary building blocks of permutation-polynomial theory. Hermite's criterion supplies a converse with an additional root condition.
\end{solution}

\begin{solution}{pb:aks}
\idea Expand with the binomial theorem: the identity holds if and only if $n$ divides all the binomial coefficients $\binom nk$,
$0<k<n$, and $a^n-a$. For a prime this is Volume~2. For a composite, one coefficient escapes, and the natural candidate is
$\binom nq$ for a prime factor $q$: its $q$-adic valuation is easy to compute.

\writeup Expand the polynomial difference:
\[
(X+a)^n-(X^n+a)
=\sum_{0<j<n}\binom nj a^{n-j}X^j+(a^n-a).
\]
If $n$ is prime, every middle binomial coefficient is divisible by $n$. Fermat's little theorem also gives $a^n\equiv a\pmod n$. Thus all coefficients of the difference vanish modulo $n$.

Now suppose $n$ is composite. Choose a prime divisor $q$ and let $q^k$ be the exact power of $q$ dividing $n$. Because $q<n$, the coefficient of $X^q$ is one of the middle coefficients and equals $\binom nq a^{n-q}$. In
\[
\binom nq=\frac{n(n-1)\cdots(n-q+1)}{q!},
\]
only the first of the $q$ consecutive numerator factors is divisible by $q$, and its valuation is $k$. The denominator has valuation $v_q(q!)=1$, so
$v_q(\binom nq)=k-1$. Since $\gcd(a,n)=1$, multiplying by $a^{n-q}$ does not add a factor of $q$. The coefficient of $X^q$ is therefore not divisible by $q^k$, hence not by $n$. The congruence fails for every composite $n$.

\tools Volume~2 (valuations, binomial coefficients and Fermat's theorem).
\leads The direct polynomial criterion is exact but not efficient: it has too many coefficients. AKS introduced a way to test a related polynomial congruence in a much smaller quotient ring.
\end{solution}

\begin{solution}{pb:wolsten}
\idea Write $\binom{2p-1}{p-1}=\prod_{k=1}^{p-1}\frac{p+k}k=\prod\bigl(1+\frac pk\bigr)$. Modulo $p^3$ only the first three elementary symmetric
functions of the $\frac pk$ matter: $1+pH_{p-1}+p^2e_2$. So we need $H_{p-1}\equiv0\pmod{p^2}$ and $e_2\equiv0\pmod p$, and both come down
to the sum of the $\frac1{k^2}$, which is a sum of squares in disguise.

\writeup All fractions below have denominators coprime to the prime $p\ge5$, so they can be interpreted modulo powers of $p$.

(a) Taking reciprocals permutes the nonzero residue classes modulo $p$. Therefore
\[
\sum_{k=1}^{p-1}\frac1{k^2}
\equiv\sum_{k=1}^{p-1}k^2
=\frac{(p-1)p(2p-1)}6
\equiv0\pmod p,
\]
since $6$ is invertible modulo $p$.

(b) Pair the term $1/k$ with $1/(p-k)$:
\[
2H_{p-1}
=\sum_{k=1}^{p-1}\left(\frac1k+\frac1{p-k}\right)
=p\sum_{k=1}^{p-1}\frac1{k(p-k)}.
\]
Inside the sum, reduce modulo $p$: $1/[k(p-k)]\equiv-1/k^2$. Part (a) makes the sum zero modulo $p$, so $2H_{p-1}$ is divisible by $p^2$. As $p$ is odd, $H_{p-1}$ is also zero modulo $p^2$.

(c) Use the product form of the binomial coefficient:
\[
\binom{2p-1}{p-1}
=\prod_{k=1}^{p-1}\left(1+\frac pk\right)
=1+pe_1+p^2e_2+O(p^3),
\]
where $e_1=H_{p-1}$ and
\[
e_2=\sum_{j<k}\frac1{jk}
=\frac12\left(H_{p-1}^2-\sum_{k=1}^{p-1}\frac1{k^2}\right).
\]
Part (b) gives $pe_1\equiv0\pmod{p^3}$; parts (a) and (b) give $e_2\equiv0\pmod p$, so $p^2e_2\equiv0\pmod{p^3}$. All higher terms already contain $p^3$. Hence
\[
\binom{2p-1}{p-1}\equiv1\pmod{p^3},
\qquad
\binom{2p}{p}\equiv2\pmod{p^3}.
\]

\tools Volume~2 (modular inverses); Volume~1 (symmetric polynomials).
\leads Wolstenholme's theorem is an early example of a binomial congruence whose modulus is much higher than the first power of a prime.
\end{solution}

\begin{solution}{pb:jacobsthal}
\idea Split $ap$ objects into $(a-1)p$ and $p$, and use Vandermonde. The terms where the split is a multiple of $p$ give the
Pascal recursion for $\binom ab$; all other terms are products of two binomial coefficients each divisible by $p$.

\writeup Proceed by induction on $a$, allowing every $0\le b\le a$. If $b=0$ or $b=a$, both sides are $1$. For $a>b\ge1$, Vandermonde's identity gives
\[
\binom{ap}{bp}
=\sum_i\binom{(a-1)p}{i}\binom p{bp-i}.
\]
A term can survive modulo $p^2$ only if at least one factor is not divisible by $p$. The boundary indices $i=bp$ and $i=(b-1)p$ give respectively
\[
\binom{(a-1)p}{bp}
\quad\text{and}\quad
\binom{(a-1)p}{(b-1)p}.
\]
For every other admissible index, $0<bp-i<p$, so $p\mid\binom p{bp-i}$. Also, $p\nmid i$ and the identity
\[
i\binom{(a-1)p}{i}
=(a-1)p\binom{(a-1)p-1}{i-1}
\]
imply $p\mid\binom{(a-1)p}{i}$. Such terms have two factors of $p$ and disappear modulo $p^2$.

Apply the induction hypothesis to the two surviving terms. Pascal's identity then yields
\[
\binom{ap}{bp}
\equiv\binom{a-1}{b}+\binom{a-1}{b-1}
=\binom ab\pmod{p^2}.
\]
This completes the induction.

\tools Volume~1 (Vandermonde and Pascal identities); Volume~2 (binomial divisibility).
\leads The mod-$p^2$ congruence is an elementary entry point to stronger congruences of binomial coefficients in $p$-adic arithmetic.
\end{solution}

\begin{solution}{pb:morley}
\idea $\binom{p-1}k=\prod_{j=1}^k\frac{p-j}j=(-1)^k\prod_j\bigl(1-\frac pj\bigr)\equiv(-1)^k(1-pH_k)\pmod{p^2}$. So the question is about the half
harmonic sum $H_{(p-1)/2}$ modulo $p$, and Problem~\ref{pb:fermatquot}(b) is exactly an alternating harmonic sum, which is
$H_{p-1}$ minus twice the even terms, that is minus $H_{(p-1)/2}$.

\writeup Put $k=(p-1)/2$. Expand the numerator of the central binomial coefficient as a product:
\[
\binom{p-1}{k}
=(-1)^k\prod_{j=1}^k\left(1-\frac pj\right)
\equiv(-1)^k(1-pH_k)\pmod{p^2},
\]
where $H_k=\sum_{j=1}^k1/j$.

The alternating harmonic sum gives a useful expression for $H_k$. Split the ordinary harmonic sum into even and odd indices:
\[
\sum_{j=1}^{p-1}\frac{(-1)^{j-1}}j
=H_{p-1}-2\sum_{j=1}^k\frac1{2j}
=H_{p-1}-H_k.
\]
Pairing $j$ with $p-j$ shows $H_{p-1}\equiv0\pmod p$, so the alternating sum is congruent to $-H_k$ modulo $p$. The Fermat-quotient identity of Problem~\ref{pb:fermatquot}(b) identifies this with $2q_p(2)$; thus
$H_k\equiv-2q_p(2)\pmod p$.

Substitute into the binomial product:
\[
\binom{p-1}k
\equiv(-1)^k(1+2pq_p(2))\pmod{p^2}.
\]
By definition, $2^{p-1}=1+pq_p(2)$. Therefore
\[
4^{p-1}=(2^{p-1})^2
\equiv1+2pq_p(2)\pmod{p^2}.
\]
The two right sides agree, proving
$\binom{p-1}{(p-1)/2}\equiv(-1)^{(p-1)/2}4^{p-1}\pmod{p^2}$.

\tools Volume~2; Problems~\ref{pb:wolsten} and~\ref{pb:fermatquot}.
\leads Morley's stronger congruence modulo $p^3$ is a natural further step. Similar central-binomial congruences occur in modern supercongruence questions.
\end{solution}

\begin{solution}{pb:glaisher}
\idea Sum the identity $k^{p-1}=1+pq_p(k)$ over $k=1,\dots,p-1$. Lerch's formula turns the right side into $p-1+pW_p$, that is
$(p-1)!+p$. The left side is a sum of powers, which Faulhaber's formula writes with Bernoulli numbers; modulo $p^2$ only the
term with $B_{p-1}$ survives, and von Staudt--Clausen controls the denominators.

\writeup The theorem about Bernoulli denominators from Volume~5 says that, for an even positive index $j$, the denominator of $B_j$ is the product of primes $\ell$ for which $\ell-1\mid j$. In particular, $B_j$ is $p$-integral for $j<p-1$, and $pB_{p-1}$ is also $p$-integral.

Use Faulhaber's formula with $n=p$ and exponent $p-1$:
\[
\sum_{k=0}^{p-1}k^{p-1}
=\sum_{j=0}^{p-1}\binom pj B_jp^{p-1-j}.
\]
The term with $j=p-1$ is $pB_{p-1}$. The term $j=0$ is $p^{p-1}$, divisible by $p^2$ for $p\ge3$. For $1\le j\le p-2$, the binomial coefficient is divisible by $p$, the displayed power of $p$ contributes at least one more, and $B_j$ has no $p$ in its denominator. Consequently
\[
\sum_{k=1}^{p-1}k^{p-1}
\equiv pB_{p-1}\pmod{p^2}.
\]

There is a second way to express the sum. The Fermat quotient satisfies
$k^{p-1}=1+pq_p(k)$, so
\[
\sum_{k=1}^{p-1}k^{p-1}
=p-1+p\sum_{k=1}^{p-1}q_p(k).
\]
Lerch's congruence from Problem~\ref{pb:fermatquot}(c) replaces the final sum by the Wilson quotient $W_p$ modulo $p$. Since $pW_p=(p-1)!+1$, the expression is congruent to $(p-1)!+p$ modulo $p^2$.

Equating the two evaluations gives
\[
(p-1)!\equiv pB_{p-1}-p\pmod{p^2},
\]
as required.

\tools Volume~5 (Faulhaber's formula and von Staudt--Clausen); Volume~2 (Wilson and Fermat quotients).
\leads This relation is a bridge between Bernoulli denominators and factorial congruences. It reduces to Wilson's theorem modulo $p$, but contains additional information modulo $p^2$.
\end{solution}

\begin{solution}{pb:ramanujan5}
\idea Two infinite products have explicit expansions: $E(x)$ (Euler) and $E(x)^3$ (Jacobi). Their product $E(x)^4$ has
coefficients we can control modulo $5$, and $E(x)^4=E(x)^5/E(x)\equiv E(x^5)P(x)$ links them to partitions.

\writeup (a) The product converges absolutely and locally uniformly on $\C^*$ since $\sum|q|^{2n-1}<\infty$, so $H$ is holomorphic on
$\C^*$ and has a Laurent expansion. Replacing $z$ by $q^2z$ shifts the factors:
\[
H(q^2z)=\prod_{n\ge1}(1+zq^{2n+1})(1+z^{-1}q^{2n-3})=H(z)\,\frac{1+z^{-1}q^{-1}}{1+zq}=\frac{H(z)}{zq}.
\]
In coefficients, $c_mq^{2m}=q^{-1}c_{m+1}$, that is $c_{m+1}=q^{2m+1}c_m$ for all $m\in\Z$. Since $(m+1)^2-m^2=2m+1$, induction upwards
and downwards from $m=0$ gives $c_m=q^{m^2}c_0$. So $H(z)=c_0(q)\theta(z,q)$ with $\theta(z,q)=\sum_mq^{m^2}z^m$.

(b) At $z=i$, $(1+iq^{2n-1})(1-iq^{2n-1})=1+q^{4n-2}$, while in $\theta(i,q)$ the terms $m$ and $-m$ cancel for odd $m$ (as
$i^m+i^{-m}=0$) and $m=2k$ gives $(-1)^kq^{4k^2}$. So
\[
\prod_{n\ge1}(1+q^{4n-2})=c_0(q)\sum_k(-1)^kq^{4k^2}.
\]
At $z=-1$ with $q^4$ in place of $q$, $\prod_n(1-q^{8n-4})^2=c_0(q^4)\sum_k(-1)^kq^{4k^2}$. The series $\sum(-1)^kq^{4k^2}=1-2q^4+\dots$ does
not vanish for small $|q|$; dividing, and using $1-q^{8n-4}=(1-q^{4n-2})(1+q^{4n-2})$,
\[
\frac{c_0(q)}{c_0(q^4)}=\frac1{\prod_n(1-q^{4n-2})\prod_n(1-q^{8n-4})}.
\]
Splitting $E(q^2)=\prod(1-q^{2n})$ according to the parity of $n$, and then $\prod(1-q^{4n})$ in the same way,
$E(q^2)=E(q^8)\prod(1-q^{8n-4})\prod(1-q^{4n-2})$. Hence $d(q)=c_0(q)E(q^2)$ satisfies $d(q)=d(q^4)$ for small $|q|$. Now
$c_0(q)=\frac1{2\pi}\int_0^{2\pi}H(e^{it})\,dt$ is holomorphic in $q$ on the unit disc and tends to $1$ as $q\to0$ (the product tends to
$1$ uniformly on $|z|=1$). So $d(q)=d(q^{4^k})\to1$ as $k\to\infty$, $d=1$ near $0$, and by analytic continuation $c_0(q)=\frac1{E(q^2)}$
on the whole disc. Multiplying $H(z)=c_0(q)\theta(z,q)$ by $E(q^2)$ gives the triple product.

(c) Take $0<x<1$, $q=x^{1/2}$ and $z=-wx^{1/2}$ with $w>0$. Then $1+zq^{2n-1}=1-wx^n$, $1+z^{-1}q^{2n-1}=1-w^{-1}x^{n-1}$ and
$q^{m^2}z^m=(-1)^mw^mx^{m(m+1)/2}$, so
\[
E(x)\prod_{n\ge1}(1-wx^n)\prod_{n\ge1}(1-w^{-1}x^{n-1})=\sum_{m\in\Z}(-1)^mw^mx^{m(m+1)/2}.
\]
The factor $n=1$ of the last product is $1-w^{-1}$. On the right, the terms $m$ and $-1-m$ have the same exponent of $x$ and
opposite signs, so the right side is $\sum_{m\ge0}(-1)^mx^{m(m+1)/2}(w^m-w^{-m-1})$. Divide both sides by $1-w^{-1}$ and let $w\to1$:
$\frac{w^m-w^{-m-1}}{1-w^{-1}}\to2m+1$ (both numerator and denominator vanish at $w=1$, with derivatives $2m+1$ and $1$), and the left
side tends to $E(x)\cdot E(x)\cdot E(x)$. The exchange of limit and sum is justified by uniform convergence for $w\in[\frac12,2]$.
Both sides of Jacobi's identity are power series in $x$ that agree on $]0,1[$, so they agree for $|x|<1$.

(d) By (c) and the pentagonal theorem,
\[
xE(x)^4=\sum_{j\in\Z}\sum_{k\ge0}(-1)^{j+k}(2k+1)x^{e(j,k)},\qquad e(j,k)=1+\frac{j(3j-1)}2+\frac{k(k+1)}2.
\]
Now
\[
8e(j,k)=(2k+1)^2+12j^2-4j+7\equiv(2k+1)^2+2(j+4)^2\pmod5,
\]
since $2(j+4)^2=2j^2+16j+32\equiv2j^2+j+2\equiv12j^2-4j+7$. The squares modulo $5$ are $0,1,4$, and twice them are $0,2,3$; a
square plus twice a square is $\equiv0$ only if both are $0$. So $5\mid e(j,k)$ forces $5\mid2k+1$, and every term contributing to a
power $x^{5m}$ has a coefficient divisible by $5$.

(e) In $\Z[[x]]$, $(1-x^k)^5\equiv1-x^{5k}\pmod5$ (the middle binomial coefficients are divisible by $5$), so $E(x)^5\equiv E(x^5)$ and
\[
xE(x)^4=\frac{xE(x)^5}{E(x)}\equiv E(x^5)\sum_{n\ge0}p(n)x^{n+1}\pmod5.
\]
Split $\sum p(n)x^{n+1}=\sum_{r=0}^4x^rA_r(x^5)$ by the residue of the exponent modulo $5$. Multiplying by $E(x^5)$ does not mix the
residues, so the part of $xE(x)^4$ with exponents divisible by $5$ is $E(x^5)A_0(x^5)$, which is $\equiv0\pmod5$ by (d). Since
$E(x^5)$ is invertible in $\Z[[x]]$, $A_0\equiv0\pmod5$. The coefficients of $A_0(x^5)$ are the $p(n)$ with $n+1\equiv0\pmod5$, that
is, the $p(5m+4)$.

\tools Volume~5, Theorem~2.12 (pentagonal number theorem); Volume~5, Chapters~2 and~11 (formal and holomorphic power series,
Laurent series); Volume~2, Chapter~3 (squares modulo $5$).
\leads This is Ramanujan's proof from 1919. The same idea with $E(x)^6=E(x)^3E(x)^3$ gives $p(7n+5)\equiv0\pmod7$, and the
congruence modulo $11$ needs more (Ramanujan's unpublished manuscript, edited by Hardy in 1921). Ramanujan also proved the
beautiful identity $\sum_np(5n+4)x^n=5\frac{E(x^5)^5}{E(x)^6}$. In 2000 Ken Ono showed that there are congruences of the form
$p(An+B)\equiv0\pmod M$ for every $M$ prime to $6$, through the theory of modular forms (Volume~6, Chapter~12, and Part~III).
\end{solution}

\begin{solution}{pb:kloosterman}
\idea Second and fourth moments of a family of exponential sums are counts of solutions of equations, which are easy. Since a
single term is at most the whole fourth moment, a bound for the moment is a bound for each sum.

\writeup (a) The substitution $x\mapsto-x$ permutes $\mathbb F_p^\times$ and changes $\bar x$ into $-\bar x$, so it turns $K(a)$ into $\overline{K(a)}$: hence $K(a)$
is real. $K(0)=\sum_{x\ne0}e(\frac{\bar x}p)=\sum_{y\ne0}e(\frac yp)=-1$. For $b\ne0$, substitute $x=by$: then $b\bar x=\bar y$ and $ax=aby$, so $K(a,b)=K(ab,1)$.

(b) Expanding, and summing over $a$ first,
\[
\sum_aK(a)\overline{K(a)}=\sum_{x,y\ne0}e\Bigl(\frac{\bar x-\bar y}p\Bigr)\sum_ae\Bigl(\frac{a(x-y)}p\Bigr)=p\sum_{x\ne0}1=p(p-1).
\]

(c) Let $M=\sum_aK(a)^4$. Compute $\Sigma=\sum_{a,b\in\mathbb F_p}|K(a,b)|^4$ in two ways. First, for $b\ne0$, (a) and the bijection $a\mapsto ab$ give
$\sum_a|K(a,b)|^4=M$; for $b=0$, $K(a,0)=\sum_{x\ne0}e(\frac{ax}p)$ is $p-1$ for $a=0$ and $-1$ otherwise. So
$\Sigma=(p-1)M+(p-1)^4+(p-1)$. Second, $|K(a,b)|^4=K(a,b)^2\,\overline{K(a,b)}^2$ expands into
$\sum_{x_1,\dots,x_4}e\bigl(\frac{a(x_1+x_2-x_3-x_4)+b(\bar x_1+\bar x_2-\bar x_3-\bar x_4)}p\bigr)$, and summing over $a$ and $b$ gives $\Sigma=p^2N$, where $N$ counts the
$(x_1,\dots,x_4)\in(\mathbb F_p^\times)^4$ with
\[
x_1+x_2=x_3+x_4,\qquad\bar x_1+\bar x_2=\bar x_3+\bar x_4.
\]
Let $s=x_1+x_2$. If $s\ne0$, then $\bar x_1+\bar x_2=\frac s{x_1x_2}$, so the equations say $x_3+x_4=s$ and $x_3x_4=x_1x_2$: $x_3$ and $x_4$ are the
roots of $X^2-sX+x_1x_2$, that is $\{x_3,x_4\}=\{x_1,x_2\}$, which gives $2$ solutions if $x_1\ne x_2$ and $1$ if $x_1=x_2$. There are $(p-1)(p-2)$
pairs $(x_1,x_2)$ with $s\ne0$, of which $p-1$ have $x_1=x_2$. If $s=0$, then $\bar x_1+\bar x_2=0$ too, and $(x_3,x_4)$ is any pair with
$x_4=-x_3$: $(p-1)^2$ solutions in all. So
\[
N=2(p-1)(p-2)-(p-1)+(p-1)^2=3(p-1)(p-2).
\]
Comparing, $(p-1)M=3p^2(p-1)(p-2)-(p-1)^4-(p-1)$, and dividing by $p-1$,
\[
M=3p^2(p-2)-(p-1)^3-1=2p^3-3p^2-3p.
\]

(d) Every term of $M$ is non-negative, so $K(a)^4\le M<2p^3$ and $|K(a)|<2^{1/4}p^{3/4}<1.19\,p^{3/4}$. The trivial bound is $p-1$, and (b) gives
only $|K(a)|<p$, no better. The fourth moment wins as soon as $1.19\,p^{3/4}<p-1$, that is for every $p\ge5$, and by a growing factor.

\tools{Volume~2, Chapter~3 (inverses modulo $p$); Volume~1, Chapter~6 (a quadratic has at most two roots in a field).}
\leads Kloosterman introduced these sums in 1926 to count representations of integers by quaternary quadratic forms, and his
exponent $\frac34$ was enough for that purpose. In 1948 André Weil deduced from his proof of the Riemann hypothesis for curves over
finite fields that $|K(a)|\le2\sqrt p$ for $a\ne0$, the optimal exponent. Writing $K(a)=2\sqrt p\cos\theta_a$, Katz
(1988) showed that the angles $\theta_a$ are distributed, as $p\to\infty$, with density $\frac2\pi\sin^2\theta$ on $[0,\pi]$: the Sato--Tate law,
which can be observed numerically.
\end{solution}

\section*{Sources and further reading}

\begin{itemize}
\item K.~Ireland and M.~Rosen, \emph{A Classical Introduction to Modern Number Theory}, 2nd ed., Springer, 1990, Chapter~3. Congruences, the Chinese remainder theorem and the structure of $\Z/n\Z$.
\item I.~Niven, H.~S.~Zuckerman and H.~L.~Montgomery, \emph{An Introduction to the Theory of Numbers}, 5th ed., Wiley, 1991, Chapter~2. A wealth of exercises on congruences, including Hensel lifting.
\item V.~Shoup, \emph{A Computational Introduction to Number Theory and Algebra}, 2nd ed., Cambridge University Press, 2009, Chapters~2 and~4. Modular arithmetic as algorithms: inverses, the Chinese remainder theorem, fast exponentiation.
\item A.~Granville, Arithmetic properties of binomial coefficients~I: binomial coefficients modulo prime powers, \emph{Canadian Mathematical Society Conference Proceedings} 20 (1997), 253--275. Lucas' theorem (Exercise~\ref{ex:lucas}) and its far-reaching generalisations.
\item W.~R.~Alford, A.~Granville and C.~Pomerance, There are infinitely many Carmichael numbers, \emph{Annals of Mathematics} 139 (1994), 703--722. The research paper that settled the question behind Exercise~\ref{ex:korselt}.
\item C.~F.~Gauss, \emph{Disquisitiones Arithmeticae}, Leipzig, 1801, Article~2, where the sign $\equiv$ is introduced; English translation by A.~A.~Clarke, Springer, 1986.
\end{itemize}
